%--------------------------------------
% Caso de Uso: Agregar Servicio a Reservación
%--------------------------------------

\begin{UseCase}{CU18}{Agregar Servicio a Reservación}{
	Permitir que un usuario autorizado agregue servicios adicionales a una reservación existente.
}
	\UCitem{Versión}{\color{Gray}1.0}
	\UCitem{Autor}{\color{Gray}Equipo de Desarrollo}
	\UCitem{Supervisa}{\color{Gray}Coordinador del Sistema}
	\UCitem{Actor}{\hyperlink{Administrador}{Administrador} o \hyperlink{Empleado}{Empleado}}
	\UCitem{Propósito}{Añadir servicios complementarios a una reservación activa.}
	\UCitem{Entradas}{Servicio seleccionado, fecha y hora, empleado asignado, notas opcionales.}
	\UCitem{Origen}{Pantalla de detalle de reservación}
	\UCitem{Salidas}{Servicio agregado a la reservación y costos actualizados.}
	\UCitem{Destino}{Pantalla de detalle de reservación y base de datos}
	\UCitem{Precondiciones}{
		\begin{itemize}
			\item El usuario debe haber iniciado sesión.
			\item La reservación debe existir y encontrarse en un estado válido.
			\item El servicio debe existir en el catálogo.
			\item La fecha y hora del servicio deben estar dentro del período de la reservación.
		\end{itemize}
	}
	\UCitem{Postcondiciones}{
		\begin{itemize}
			\item El servicio queda asociado a la reservación.
			\item El costo total de la reservación se actualiza.
			\item El servicio queda programado en el calendario.
			\item El empleado queda asignado al servicio.
		\end{itemize}
	}
	\UCitem{Errores}{
		Servicio no disponible, conflicto de horarios, empleado no disponible, reservación en estado inválido.
	}
	\UCitem{Tipo}{Caso de uso secundario}
	\UCitem{Observaciones}{
		\begin{itemize}
			\item Cada servicio puede requerir un tipo específico de empleado.
			\item Se valida la disponibilidad del empleado antes de confirmar.
			\item El costo del servicio puede variar según características de la mascota.
		\end{itemize}
	}
\end{UseCase}

%--------------------------------------
\begin{UCtrayectoria}
	\UCpaso[\UCactor] Accede a la pantalla de detalle de una reservación.
	\UCpaso[\UCactor] Selecciona la opción para agregar un servicio.
	\UCpaso El sistema muestra el catálogo de servicios disponibles.
	\UCpaso[\UCactor] Selecciona un servicio.
	\UCpaso El sistema muestra el formulario de configuración del servicio.
	\UCpaso[\UCactor] Selecciona la mascota que recibirá el servicio.
	\UCpaso[\UCactor] Indica la fecha y hora del servicio.
	\UCpaso El sistema valida que la fecha y hora estén dentro del período de la reservación.
	\UCpaso El sistema verifica la disponibilidad del servicio.
	\UCpaso[\UCactor] Selecciona un empleado o permite asignación automática.
	\UCpaso El sistema verifica la disponibilidad del empleado.
	\UCpaso[\UCactor] Ingresa notas opcionales para el servicio.
	\UCpaso[\UCactor] Confirma la adición del servicio.
	\UCpaso El sistema calcula el costo del servicio.
	\UCpaso El sistema registra el servicio en la reservación.
	\UCpaso El sistema actualiza el costo total de la reservación.
	\UCpaso El sistema agenda el servicio en el calendario.
	\UCpaso El sistema muestra un mensaje indicando que el servicio fue agregado correctamente.
	\UCpaso El sistema actualiza la pantalla de detalle de la reservación.
\end{UCtrayectoria}

%--------------------------------------
\begin{UCtrayectoriaA}{A}{Fecha u hora fuera del período de la reservación}
	\UCpaso El sistema informa que el servicio debe programarse dentro del período de la reservación.
	\UCpaso El flujo continúa desde la selección de fecha y hora.
\end{UCtrayectoriaA}

%--------------------------------------
\begin{UCtrayectoriaA}{B}{Empleado no disponible}
	\UCpaso El sistema informa que el empleado no se encuentra disponible en la fecha y hora seleccionadas.
	\UCpaso El usuario puede modificar la fecha, seleccionar otro empleado o cancelar la operación.
\end{UCtrayectoriaA}

%--------------------------------------
\subsection{Puntos de extensión}
\UCExtenssionPoint{
	El usuario desea agregar varios servicios a la misma reservación.
}{
	Del paso 4.
}{
	Agregar Múltiples Servicios.
}

%--------------------------------------
% Fin del Caso de Uso
